
% Conclusiones de la práctica experimental GA-SUM-01.
% Cada integrante debe redactar su propia conclusión de mínimo 80 palabras.

\section*{Conclusiones}
\addcontentsline{toc}{section}{Conclusiones}

\begin{enumerate}

\item \textbf{Conclusión de Luis Yépez}

Durante esta práctica comprendí que para desarrollar un sistema no basta con pensar en las funciones que tendrá. Primero debemos identificar las necesidades de los usuarios y conocer el contexto donde se utilizará. En el proyecto UteqFit-AI analizamos los tipos de requisitos, sus propiedades y las ocho perspectivas de contexto que presenta Pohl (2025)~\cite{pohl2025}. Esto nos permitió identificar aspectos importantes como el control de asistencias, la protección de los datos de los estudiantes y las posibles rutinas con inteligencia artificial. También aprendí que los requisitos deben ser claros, verificables y revisados con los responsables del gimnasio~\cite{ieee29148}. Considero que este trabajo nos ayudará a organizar mejor el proyecto y evitar problemas durante su desarrollo.

% Alan y Anderson agregarán aquí las otras cuatro conclusiones.

\end{enumerate}
